% Реферат
\chapter*{РЕФЕРАТ}
\addcontentsline{toc}{chapter}{РЕФЕРАТ}

Расчетно-пояснительная записка содержит \ruquantity{\getpagerefnumber{LastPage}}{страницу}{страницы}{страниц}, \ruquantity{\totvalue{figures}}{рисунок}{рисунка}{рисунков}, \ruquantity{\totvalue{tables}}{таблицу}{таблицы}{таблиц}, 1 приложение, \ruquantity{\totvalue{references}}{источник}{источника}{источников}.

Ключевые слова: медицинская клиника, база данных, PostgreSQL, Redis, Go, Gin, React, TypeScript, UUID, триггеры, хранимые процедуры, индексы, конкурентный доступ.

Объектом разработки является информационная система медицинской клиники, обеспечивающая регистрацию пациентов, ведение расписания врачей, запись на прием, проведение приема, формирование счетов и учет оплаты услуг.

Цель работы --- спроектировать и реализовать базу данных и приложение доступа к ней для автоматизации деятельности медицинской клиники, а также исследовать вопросы согласованности данных и производительности ключевых операций.

В ходе выполнения работы проведены анализ предметной области медицинской клиники, формализация задачи, проектирование структуры базы данных, ролевой модели и прикладного интерфейса доступа. Реализованы схема базы данных PostgreSQL, хранимые процедуры, триггеры, сценарии тестирования и кэширование расписания врачей и прайс-листа услуг с использованием Redis. Разработан backend на Go и Gin с архитектурой Repository $\rightarrow$ Service $\rightarrow$ Handler и подготовлен frontend на React и TypeScript. Выполнены проверки валидации, HTTP-уровня, кэширования, конкурентного доступа и производительности ключевых операций.

\clearpage
